\section{Parameter choice and completion of the proof}\label{sec:completion}
We finish Proposition~\ref{prop:normalized}.  Throughout this section all
collections of summation cells are finite and lie in the original $N$
depths.  Thus the defining constant $C_N(q)$ is finite, and every
application of it in Section~\ref{sec:sides} concerns precisely this
finite constant.

\subsection{One attempt and its restarts}
Fix one of the normalized base segments from Section~\ref{sec:base},
with root $R$, a tuple of increment levels, and one active pattern.
For the noninitial levels, its remaining quantity is
\[
 \sum_{I\text{ active}}|I|\min\{\eta,|H_I(z)|\},
 \qquad \eta=e^{-k^\theta},
\]
where $k$ is the later accuracy log and $\theta>0$ is fixed.
This is Equation~\eqref{eq:src-11}.
We explain explicitly how the stopping constructions preserve this
quantity.

Start an attempt at $R$.  Construct the main catalogs, color histories,
and ball filters, making the small marking-failure cuts from
Section~\ref{sec:localization}.  The objects needed on an unused branch
can be continued by inheritance and local refitting.  Next construct the
fixed first-match mode expansions and the linear side states of
Section~\ref{sec:linearization}.  Choose the branch offsets by the small
integrated error bound over all the pre-existing actual uses.  Finally
close the helper catalogs on the entire side states and the specified
histories.  A helper base is one such state, not an individual Fourier
mode or a choice of heights and amplitude bins.

All subsequent size and counting stops are imposed on these fixed
objects.  A stop removes original summands from the attempt; it does not
truncate or redefine a previously fixed stack function.  Partial path
counts and ancestor suprema are accumulated in increasing depth, and a
cell is stopped at their first threshold crossing.  Hence the retained
summands see only retained partial counts.  For the hull estimates of
Section~\ref{sec:sides}, include all potentially needed ancestors within
a block before imposing the common cuts.  The maximal functions remain
defined on the full retained hull cells, including portions below a cut.
Their joint caps are consequently valid on the sparse remainders used
in the proof of Equation~\eqref{eq:src-19}.
Large-frozen-complement flags are charged as powered errors within the
current attempt; only a first crossing of their accumulated power-cost
budget produces a restart root.

The integral bounds proved in the preceding sections permit each class
of stops to have arbitrarily small fixed relative root length by
increasing its threshold by a fixed factor.  The number of configurations
requiring the strong caps is a small overhead, which can be included in
those thresholds.  Polynomially many configurations in $P,k$ occur only
in the estimates already permitted a polynomial cost.  Choose these
thresholds together so that the union of first stopped cells has total
length at most, say, $|R|/4$.

At every stopped cell send the original base summands there and below
to a fresh attempt.  The original basic arguments and their inherited
lists remain the same.  Their normalized sizes hold on the new root;
the rolling basic projection bounds and the other local hypotheses hold
there as well.  Thus each attempt has the same estimates, uniformly in
its root and in the remaining number of depths.  The total length of
all roots in the resulting recursion is at most
\[
 |R|\sum_{n\ge0}4^{-n}=\tfrac43|R|.
\]
The construction terminates on a finite depth tree.  Equivalently, one
can prove the required bound by induction on the remaining depth, with
the displayed geometric root-length estimate.

\subsection{The estimate on retained cells}
Proposition~\ref{prop:localization} first gives the group decomposition
in Equation~\eqref{eq:src-14}.  In each coarse role the decomposition of
Section~\ref{sec:linearization} is a core plus a rolling and a linear
side, with a separately estimated physical error.  The basic rolling
$U$ role, if present, is always a side.

Terms with at least three core roles are estimated by
Proposition~\ref{prop:core-estimate}.  Terms with at least three side
factors have three square-size budgets, so they are summable at every
fixed power $p>2/3$ sufficiently close to one, with a polynomial cost in
$P,k$.  In a term with exactly two sides, replace its two core complements
by their whole current localized arguments.  The correction terms have
at least three sides.  Apply Lemma~\ref{lem:transfer} to the two whole
complements.  After transfer, replacing their basic coarse parts by the
full $Y$ arguments adds a third basic side; the remaining term has the
full complementary pair for which Section~\ref{sec:base} proves
smallness.  This is the pair-small main term in
Proposition~\ref{prop:side-estimate}.

The near-band defects are retained as their coupled band sums and
estimated at power one by Equation~\eqref{eq:src-19}.  Numerical errors
are also estimated at power one.  For all the other terms we use
\[
 \min\{\eta,a\}\le\eta^{1-p}a^p\qquad(a\ge0).
\]
The minimum is first applied to the sum at the original cell and then
split by subadditivity.  In particular no separate pair-smallness
assertion for an individual group is being used.  Put $p=1-\varepsilon$,
where $\varepsilon>0$ is fixed small enough for both the three-side
estimate and the interpolation in Section~\ref{sec:counting}.

The first term below collects the estimates taken at power
$p=1-\varepsilon$ and the smaller pair-small main term; the second is the
coupled near-band estimate at power one.  Combining
Equations~\eqref{eq:src-19} and~\eqref{eq:src-37} with the other side
estimates, and increasing one fixed exponent $C_4$ to include all powered
costs, gives on retained cells
\[
 C_q\left[
 \eta^{\varepsilon}P^{C_4}
 +(1+C_N)P^{-1/2}G_0(K+2)^{C_{19}(q-2)}
 \right]|R|+O(e^{-k}|R|).
\]
Fixed polynomial factors in $k$ can be retained here and absorbed into
the first term's power of $P$ once the parameters below are fixed.  The
main term after transfer, of size $\eta P^C|R|$, is also covered after
increasing $C_4$ and decreasing $\varepsilon$ if necessary.  The constant
$C_{19}$ is absolute.  Restarting as above changes this bound only by a
fixed factor.

\subsection{Order of the parameters}
For clarity, the two types of complexity used in the proof are listed
separately:
\begin{center}
\begin{tabular}{p{.28\linewidth}p{.31\linewidth}p{.30\linewidth}}
\toprule
Quantity & Allowed size & Role in the argument\\
\midrule
$D_{\mathrm{base}},D_{\mathrm m}$
 & $\log D\le C_q k^C$
 & Main columns, colors, helpers and their projection products.\\[3pt]
$D_*$ and the shift heights
 & $\log D_*\le C_q(1+k+P)^C$
 & Fixed linear modes, Fourier truncations and graph tests.\\[3pt]
$K$
 & $K\le C_q(1+k+P)^{C_5}$
 & Depth-block length after the mode and gap choices.\\[3pt]
$G_0$
 & $\exp(C_q(1+\log(2k))^C)$
 & Strong energy and stack overhead in Equations~\eqref{eq:src-17}--\eqref{eq:src-19}.\\
\bottomrule
\end{tabular}
\end{center}
Every exponent denoted by $C,C_4,C_5,C_{19}$ in this hierarchy has an
absolute upper bound before $q$ is chosen.  Constants such as $C_q$ may
depend on that later fixed choice.

Here is why the ordering is not circular.  Detection changes a power of
an accuracy log by an absolute amount; the dependence on $q$ from
clipping multiplies that log by a $q$-dependent constant.  Thus choose
the stagger exponent $A$ in Section~\ref{sec:base} and the small quality
exponent $\theta$ independently of $q$.  The main/helper catalogs have
polynomially many ranks and multiplicative count increments
$\exp(k^{C_2})$.  Their logarithms remain polynomial in $k$.
The number of masks and histories enters this recurrence logarithmically,
including the $P$ pad choices, because $\log P=O(k)$.  The tolerance in
Equation~\eqref{eq:src-13} includes $n_{r-1}^{-1/2}$, so summing a lower-rank
source does not require its full count as an additional error loss.
Choose the mismatch accuracy and the edge height after these count
bounds.  Color recourse bounds are uniform in that edge height.

The short-path construction then produces the larger mode budget.
Its Fourier cutoffs precede the height range, shift radii and horizontal
gap, and $K$ is chosen last among these quantities.  All their logarithms
are still polynomial in $P,k$.  The strong side estimate compresses only
main projection families; mode projections have already been incorporated
in Bessel rows.  The bad-density argument removes the large height
factor from its chain count.  Thus $G_0$ has the smaller size in the table,
whereas the top counts and core estimates may pay powers of $P$.
No later mode count is fed back into main/helper closure.

Now set $P=\lceil\eta^{-\alpha}\rceil$ for a fixed $\alpha>0$ to be
chosen.  For every such fixed choice and sufficiently large $k$,
$P$ dominates each fixed power of $k$ and $G_0=P^{o(1)}$.
First choose $2<q<3$ so close to two that
\[
 C_{19}C_5(q-2)<\tfrac18.
\]
After the finite overheads are combined, increasing the starting log
makes them at most $P^{1/8}$.  Hence
\[
 P^{-1/2}G_0(K+2)^{C_{19}(q-2)}\le C_qP^{-1/4}.
\]
Next choose $\alpha>0$ small enough that
$\alpha C_4<\varepsilon/2$, and impose the finitely many analogous
upper bounds coming from the other powered side costs.  Then
\[
 \eta^{\varepsilon}P^{C_4}\lesssim\eta^{\varepsilon/2},
 \qquad P^{-1/4}\le\eta^{\alpha/4}.
\]
Finally increase the starting log so that all earlier accuracy and
small-overhead requirements hold.  Since $\theta<1$, numerical errors
$e^{-k}$ are smaller than the resulting quality gain for large $k$.
We have therefore proved, for some fixed $c_*>0$, the bound
\[
 \sum_{I\text{ active}}|I|\min\{\eta,|H_I(z)|\}
 \le C_q\eta^{c_*}(1+C_N)|R|
\]
on every base segment and pattern.

\subsection{Summing the base expansion}
For later level $t$, there are $O(t)$ choices of the earlier increment
level and only a bounded number of slot patterns.  Each tuple uses
intersections of a bounded number of inheritance segments.  Their roots
have uniformly bounded total length by the reset packing in
Section~\ref{sec:base}; the further attempts just considered preserve
this property.  Thus the total coefficient of $1+C_N$ from the
noninitial tuples is bounded by a fixed multiple of
\[
 \sum_{t\ge12}t\exp(-c_*k_t^\theta),
 \qquad k_{t+1}=k_t^A.
\]
This series converges and tends to zero as $k_1$ tends to infinity.
The absolute errors $O(e^{-k_t})$ are summable with the same multiplicity.
The finitely many initial levels are bounded by the active-scale count
in Equation~\eqref{eq:src-10}; their cost is a finite constant depending
on the chosen starting logs but independent of $N$.

Given the small $c$ in Proposition~\ref{prop:normalized}, take $k_1$
large enough that the preceding coefficient of $C_N$ is at most $c$.
All remaining constants are included in $A$.  The base expansion is
first taken at a finite accuracy cutoff; on a fixed finite tree its
residual terms tend to zero by the staggered pair comparisons of
Section~\ref{sec:base}.  The uniform bounds just proved permit passage
to that limit.  We obtain the assertion of
Proposition~\ref{prop:normalized}.
Lemma~\ref{lem:outer} now proves Theorems~\ref{thm:local}
and~\ref{thm:main}.
